\subsection{Angular stability for unequal push-outs}
\label{round26:sec:root-angular}

Let $e_1,\ldots,e_4$ be the coordinate basis and let
\[
 \mathcal U=\{(\pm e_i\pm e_j)/\sqrt2:1\le i<j\le4\}.
\]
Write $u_1,\ldots,u_{24}$ for these unit vectors, and put
\[
 Q=\{x:\langle x,u_i\rangle\le1\text{ for all }i\},
 \qquad K_0=\operatorname{conv}\{2u_i\}.
\]
The eight axis vertices of $Q$ have the form $\pm\sqrt2e_j$.
The remaining sixteen vertices have the form $s/\sqrt2$ with $s\in\{-1,1\}^4$. Thus
$Q\subset B(\sqrt2)$ and $B(\sqrt2)\subset K_0$.
For completeness, the volume of $Q$ is $8$. In a sector with
$0\le x_4\le x_3\le x_2\le x_1$ the only remaining inequality
is $x_1+x_2\le\sqrt2$. Its volume is
\[
 \int_0^{1/\sqrt2}\frac{x_2^2}{2}(\sqrt2-2x_2)\,dx_2
 =\frac1{48}.
\]
Multiplication by $2^4\cdot4!=384$ gives $8$.

Let $\eta_i\in[0,1/10]$, let $w_i$ be unit vectors, and define
\[
 y_i=2(1+\eta_i)w_i,\quad
 S=\sum_i\eta_i,\quad \epsilon=\max_i|w_i-u_i|,
\]
\[
 V_w=\{x:\langle x,w_i\rangle\le1+\eta_i\text{ for all }i\},
 \quad K_w=\operatorname{conv}\{2w_i/(1+\eta_i)\}.
\]
The directions can be rotated and labelled before computing $\epsilon$.\footnote{The angular error is measured after one common rotation and one root labelling. The condition below uses the total radial displacement; it imposes no positive lower bound on each individual displacement.}

\begin{thm}\label{round26:thm:unequal-angular}
If $\epsilon\le S/500$, then
\[
 \operatorname{vol}(V_w\cap K_w)\ge8+\frac S4.
\]
If $0,y_1,\ldots,y_{24}$ are the centres within $\sqrt6$ of the
origin in a packing of unit balls, the actual Voronoi cell has the
same lower bound. Equality with $8$ forces $S=0$, $w_i=u_i$, and
the regular twenty-four-cell.
\end{thm}

This theorem permits some or all but one of the $\eta_i$ to vanish.
It therefore covers a different sector from the earlier bound that
requires every push-out to be at least half the largest push-out.

\begin{proof}
First put $w_i=u_i$, and denote the resulting bodies by $V_\eta$
and $K_\eta$. We prove
\begin{equation}\label{round26:eq:pure-radial-bound}
 \operatorname{vol}(V_\eta\cap K_\eta)\ge8+\frac{3S}{8}.
\end{equation}

Fix $u_i=(e_1+e_2)/\sqrt2$ and use the orthonormal coordinates
$t=(x_1+x_2)/\sqrt2$, $z=(x_1-x_2)/\sqrt2$, $x_3,x_4$.
When every other support number is one, the slice at $t=1+s$,
$0\le s\le1$, is
\[
 |z|+\sqrt2\max(|x_3|,|x_4|)\le1-s.
\]
All other root inequalities follow from this one. Its three-volume
is $\frac43(1-s)^3$, as follows by integrating the square cross
sections in the $z$ direction. The additional cap formed by moving
only this facet out by $\eta_i$ therefore has volume
$\frac13(1-(1-\eta_i)^4)$. The caps formed by moving distinct
single facets have disjoint interiors outside $Q$: a point of the
$i$-th cap violates only the $i$-th original inequality. Every such
cap lies in $V_\eta$. Hence
\begin{equation}\label{round26:eq:radial-gain}
 \operatorname{vol}(V_\eta)
 \ge8+\frac13\sum_i\bigl(1-(1-\eta_i)^4\bigr).
\end{equation}

We next bound what $K_\eta$ removes. Put $h=1+\max_i\eta_i\le11/10$.
The origin belongs to $K_\eta$, since it lies on the segment joining the two vertices on any pair of opposite root rays.
Then $V_\eta\subset hQ$ and $K_\eta\supset h^{-1}K_0$.
Each facet of $K_0$ corresponds to a vertex $v$ of $Q$ and has
six vertices $2u_i$. The root symmetries are transitive on these
facets. Indeed signed permutations and the orthogonal Hadamard
matrix divided by $2$ preserve $\mathcal U$ and send every vertex
direction of $Q$ to a coordinate direction.

Take $x\in V_\eta\setminus K_\eta$. Since it lies outside
$h^{-1}K_0$, it violates one of its facet inequalities. Use the
stated symmetry to make that inequality $x_1>\sqrt2/h$.
As $x\in hQ$, one has
\[
 |x_j|\le\sqrt2h-x_1,\qquad
 \frac{|x_j|}{x_1}<h^2-1\le\frac{21}{100}\quad(2\le j\le4).
\]
Let $t_v$ be the largest $\eta_i$ among the six roots through
$v=\sqrt2e_1$. The octahedron at
$x_1=\sqrt2/(1+t_v)$ with transverse $\ell^1$ radius
$\sqrt2/(1+t_v)$ belongs to $K_\eta$: all its vertices lie on
segments from zero to the corresponding $2u_i/(1+\eta_i)$.
The ray through $x$ meets this octahedron, because the sum of
its three transverse coordinate ratios is below $63/100<1$.
Since $x$ lies outside $K_\eta$, it follows that
$x_1>\sqrt2/(1+t_v)$.
The six corresponding inequalities of $V_\eta$ also give
$|x_j|\le\sqrt2(1+t_v)-x_1$. Therefore the portion assigned to
this facet has volume at most
\[
 \int_{\sqrt2/(1+t_v)}^{\sqrt2(1+t_v)}
      8\bigl(\sqrt2(1+t_v)-x_1\bigr)^3\,dx_1
 \le128t_v^4.
\]
There are twenty-four facets, and each root belongs to six of
them. Taking a union bound gives
\begin{equation}\label{round26:eq:radial-loss}
 \operatorname{vol}(V_\eta\setminus K_\eta)
 \le768\sum_i\eta_i^4.
\end{equation}
For $0<t\le1/10$, put
\[
 q(t)=\frac{1-(1-t)^4}{3t}-768t^3
 =\frac43-2t+\frac43t^2-\frac{2305}{3}t^3.
\]
Its derivative is negative on this interval, since
$q'(t)=-2+\frac83t-2305t^2\le-26/15$.
Moreover $q(1/10)=227/600>3/8$.
Combining \eqref{round26:eq:radial-gain} and \eqref{round26:eq:radial-loss}
proves \eqref{round26:eq:pure-radial-bound}.

We now allow the $w_i$ to vary. Put $a=1-8\epsilon/5$.
Since $S\le12/5$ and $\epsilon\le S/500$, $a$ is positive.
If $x\in V_\eta$, then $|x|\le\sqrt2h<8/5$, and
\[
 \langle ax,w_i\rangle
 \le a(1+\eta_i+8\epsilon/5)\le1+\eta_i.
\]
Thus $aV_\eta\subset V_w$.
For any test vector $z$, the support functions obey
\[
 h_{K_w}(z)\ge h_{K_\eta}(z)-2\epsilon|z|.
\]
Also $h_{K_\eta}(z)\ge(\sqrt2/h)|z|$. It follows that
\[
 h_{K_w}(z)\ge(1-\sqrt2h\epsilon)h_{K_\eta}(z)
 \ge a h_{K_\eta}(z).
\]
The comparison of support functions gives $aK_\eta\subset K_w$.
Consequently \eqref{round26:eq:pure-radial-bound} yields
\begin{align*}
 \operatorname{vol}(V_w\cap K_w)
 &\ge a^4(8+3S/8)\\
 &\ge 8+3S/8-\frac{32}{5}\epsilon(8+3S/8)\\
 &\ge 8+\left(\frac38-\frac{1424}{12500}\right)S
 \ge8+S/4.
\end{align*}
Here $a^4\ge1-32\epsilon/5$ is Bernoulli's inequality, and
$8+3S/8\le89/10$.

Finally let $z$ be any packing centre outside the selected set.
The packing inequalities give
$2\langle y_i,z\rangle\le|y_i|^2+|z|^2-4$ and $|y_i|,|z|\ge2$.
Hence
\[
 2\left\langle\frac{4y_i}{|y_i|^2},z\right\rangle
 \le |z|^2,
\]
because the difference after multiplication by $|y_i|^2$ is
$(|y_i|^2-4)(|z|^2-4)\ge0$.
Convexity shows that the exterior centres do not cut $K_w$.
Thus the actual cell and $V_w$ agree there. If $S>0$ the volume
bound is strict above $8$. If $S=0$, the hypothesis gives
$\epsilon=0$, and $V_w=Q\subset K_w=K_0$, giving the claimed
cell and equality statement.
\end{proof}

\begin{example}[A noncomparable family]
Take $\eta_1=2/25$ and $\eta_i=0$ for $i\ne1$. Write
$u_1=(e_1+e_2)/\sqrt2$, $v=(e_1-e_2)/\sqrt2$, and put
\[
 w_1=\cos(10^{-4})u_1+\sin(10^{-4})v,\qquad w_i=u_i\ (i\ne1).
\]
Here $\epsilon\le10^{-4}<S/500$, so the theorem gives
$\operatorname{vol}(V_w\cap K_w)\ge401/50$.
These centres form a packing. All pairs avoiding $y_1$ are root pairs.
For the others, with $r=54/25$, the inequality
$\langle w_1,u_j\rangle\le1/2+\epsilon$ gives
\[
 |y_1-y_j|^2\ge4+r(r-2-4\epsilon)>4.
\]
The example has twenty-three zero push-outs and one positive push-out.
\end{example}

The angular hypothesis $\epsilon\le S/500$ is an explicit
restriction. No argument here puts every configuration with
$T(Y)\le8$ into this sector, and the twenty-five-, twenty-six-
and remaining twenty-seven-centre cases are not treated.


The support-function comparison used here is the standard inclusion criterion for convex bodies, derived directly in the proof; see also~\cite{Round26Schneider}.

\begin{figure}[tb]
\centering
\begin{tikzpicture}[x=2.8cm,y=6.5cm,>=Latex,font=\small]
\fill[blue!12] (0,0)--(2.4,.48)--(2.4,0)--cycle;
\draw[->] (0,0)--(2.6,0) node[right] {$S$};
\draw[->] (0,0)--(0,.54) node[above] {$\epsilon$};
\draw[blue!70!black,thick] (0,0)--(2.4,.48);
\node[anchor=west,fill=white] at (.9,.37) {$\epsilon=S/500$};
\node[align=center] at (1.5,.13) {proved sector\\$\operatorname{vol}\ge8+S/4$};
\fill[red!70!black] (.08,.01) circle[x radius=.025,y radius=.01];
\draw[red!70!black] (.10,.02)--(.48,.10) node[above,align=left] {one displaced\\neighbour};
\draw (2.4,0)--(2.4,-.008) node[below] {$12/5$};
\end{tikzpicture}
\caption{The angular stability theorem. The blue sector permits unequal radial displacements, including zero displacements. Configurations outside this sector require additional estimates.}
\label{round26:fig:d4-sector}
\end{figure}